\section{Thiết kế Rule Evaluation Engine}

Rule Evaluation Engine là thành phần quyết định một Feature Flag bật hay tắt cho từng request. Mục này là nơi duy nhất trình bày chi tiết thứ tự đánh giá, cách dùng context, stable hashing, reason code và fallback; các mục dữ liệu/API chỉ tham chiếu đến contract này ở mức tổng quan.

\subsection{Evaluation request và context}

Evaluation request gồm khóa project, environment, flag và một evaluation context do ứng dụng khách cung cấp. Context chỉ là đầu vào cho rollout/targeting, không thay thế cơ chế xác thực hay phân quyền nghiệp vụ.

\begin{lstlisting}[
    caption={Cấu trúc evaluation request},
    label={lst:evaluation-request}
]
{
  "projectKey": "ecommerce-demo",
  "environmentKey": "production",
  "flagKey": "new-checkout",
  "context": {
    "targetingKey": "user-123",
    "userId": "user-123",
    "role": "BETA_CUSTOMER",
    "attributes": {
      "organizationId": "org-01"
    }
  }
}
\end{lstlisting}

\subsubsection{Project key, environment key và flag key}

Ba khóa này xác định chính xác cấu hình cần evaluate. Nếu project hoặc flag không tồn tại, engine trả kết quả tắt an toàn với reason \texttt{NOT\_FOUND}.

\subsubsection{Targeting key}

\texttt{targetingKey} là định danh ổn định, không chứa PII, dùng cho percentage rollout. Khi thiếu trường này, rule percentage không đủ điều kiện để match.

\subsubsection{User ID, role và organization ID}

\texttt{userId} phục vụ allowlist; \texttt{role} phục vụ role targeting; \texttt{attributes.organizationId} được giữ như dữ liệu mở rộng cho demo và hướng phát triển sau này.

\subsubsection{Thuộc tính mở rộng}

\texttt{attributes} cho phép bổ sung dữ liệu mà không đổi contract chính. Engine hiện chỉ dùng một phần nhỏ, các rule theo thuộc tính nâng cao được để ở hướng phát triển.

\subsubsection{Validation evaluation request}

Backend kiểm tra kiểu dữ liệu và định dạng bắt buộc trước khi gọi engine. Request sai định dạng bị từ chối ở biên API, còn request hợp lệ nhưng thiếu dữ liệu cho một rule sẽ được engine xử lý theo hướng bỏ qua rule đó.

\subsection{Evaluation snapshot}

Snapshot là ảnh chụp cấu hình cần thiết tại thời điểm evaluate, gồm trạng thái flag, config theo environment, group config và danh sách rule đã chuẩn hóa.

\subsubsection{Khái niệm evaluation snapshot}

Snapshot giúp engine không phụ thuộc trực tiếp vào ORM hoặc controller. Cùng một snapshot và cùng một context phải luôn tạo ra cùng một kết quả.

\subsubsection{Dữ liệu cấu hình cần truy xuất}

Snapshot chứa lifecycle flag, trạng thái enabled/disabled, serving mode, kill switch, group kill switch và rule parameters. Dữ liệu audit/history không nằm trong snapshot runtime.

\subsubsection{Sắp xếp rule trước khi evaluation}

Repository chuẩn bị rule theo loại và priority để engine chỉ tập trung vào logic quyết định. Cách sắp xếp này làm kết quả ổn định giữa các lần gọi.

\subsubsection{Tính nhất quán của snapshot}

Snapshot được lấy từ database hoặc cache đã invalidation sau mutation. Nếu cache miss, backend nạp lại snapshot từ nguồn dữ liệu chính.

\subsection{Thứ tự ưu tiên khi đánh giá}

Thứ tự ưu tiên là hợp đồng quan trọng nhất của engine. Các điều kiện an toàn được xét trước rule rollout để tránh việc một flag bị tắt nhưng vẫn vô tình được bật bởi rule phía sau.

\begin{lstlisting}[
    caption={Thứ tự ưu tiên của Rule Evaluation Engine},
    label={lst:evaluation-precedence}
]
ARCHIVED_FLAG
FLAG_DISABLED
GROUP_KILL_SWITCH
KILL_SWITCH
GLOBAL_ON
USER_ALLOWLIST
ROLE_TARGETING
PERCENTAGE_ROLLOUT
DEFAULT_OFF
\end{lstlisting}

\begin{figure}[H]
    \centering
    \includegraphics[width=0.90\textwidth]{\detokenize{Images/mermaid/Thứ tự ưu tiên khi đánh giá.png}}
    \caption{Flowchart thứ tự ưu tiên của Rule Evaluation Engine}
    \label{fig:evaluation-precedence-flowchart}
\end{figure}

\subsubsection{Flag đã được archive}


Flag archived không tham gia phục vụ runtime, vì vậy engine trả tắt trước mọi cấu hình khác.

\subsubsection{Flag bị disabled}

Nếu config theo environment đang disabled, engine trả tắt ngay cả khi còn rule đang enabled.

\subsubsection{Group kill switch}

Khi group kill switch bật cho environment hiện tại, toàn bộ flag thuộc group bị tắt sau khi đã qua kiểm tra archived/disabled.

\subsubsection{Flag kill switch}

Flag-level kill switch là lớp dừng khẩn cấp cho một flag riêng lẻ.

\subsubsection{Global on}

Nếu config enabled và serving mode là global on, engine bật flag cho mọi context hợp lệ mà không cần xét rule targeting.

\subsubsection{User allowlist}

Allowlist được xét trước các rule mở rộng phạm vi vì đây là điều kiện cụ thể nhất.

\subsubsection{Role targeting}

Role targeting được xét sau allowlist để hỗ trợ phát hành theo nhóm người dùng.

\subsubsection{Percentage rollout}

Percentage rollout là lớp cuối trong nhóm rule vì nó mở rộng phát hành theo tỷ lệ.

\subsubsection{Default off}

Nếu không có điều kiện nào match, engine trả tắt an toàn.

\subsection{Đánh giá trạng thái toàn cục}

Trạng thái lifecycle, enabled/disabled, serving mode và kill switch quyết định các nhánh đầu của engine. Bảng~\ref{tab:status-serving-mode-interaction} tóm tắt các trường hợp chính.

\begin{table}[H]
    \centering
    \caption{Tương tác giữa trạng thái và serving mode}
    \label{tab:status-serving-mode-interaction}
    \begin{tabular}{p{0.22\textwidth} p{0.22\textwidth} p{0.22\textwidth} p{0.24\textwidth}}
        \toprule
        \textbf{Lifecycle/config} & \textbf{Kill switch} & \textbf{Serving mode} & \textbf{Kết quả ưu tiên} \\
        \midrule
        Archived & Bất kỳ & Bất kỳ & Tắt, \texttt{FLAG\_ARCHIVED} \\
        Disabled & Bất kỳ & Bất kỳ & Tắt, \texttt{FLAG\_DISABLED} \\
        Enabled & Group on & Bất kỳ & Tắt, \texttt{GROUP\_KILL\_SWITCH} \\
        Enabled & Flag on & Bất kỳ & Tắt, \texttt{FLAG\_KILL\_SWITCH} \\
        Enabled & Off & Global on & Bật, \texttt{GLOBAL\_ON} \\
        Enabled & Off & Targeted & Xét rule theo precedence \\
        \bottomrule
    \end{tabular}
\end{table}

\subsection{Targeting theo người dùng}

User targeting dựa trên \texttt{context.userId}. Rule chỉ match khi user ID trong context thuộc danh sách đã cấu hình.

\subsubsection{User allowlist}

Allowlist phù hợp cho kiểm thử nội bộ hoặc mở tính năng cho một nhóm nhỏ đã biết trước.

\subsubsection{So khớp user ID}

Việc so khớp dùng giá trị chính xác, không suy diễn từ targeting key để tránh nhầm giữa định danh rollout và định danh nghiệp vụ.

\subsubsection{Tách userId và targetingKey}

\texttt{userId} dùng cho allowlist; \texttt{targetingKey} dùng cho hashing. Tách hai trường giúp hệ thống dùng một khóa rollout không chứa PII.

\subsubsection{Xử lý khi thiếu định danh người dùng}

Nếu thiếu \texttt{userId}, allowlist không match nhưng engine vẫn có thể xét role hoặc percentage nếu context còn đủ dữ liệu cho các rule đó.

\subsection{Targeting theo vai trò}

Role targeting dựa trên vai trò do context cung cấp và phù hợp với demo phân nhóm người dùng.

\subsubsection{Cấu trúc role trong evaluation context}

Role có thể là một chuỗi hoặc danh sách chuỗi tùy contract triển khai; backend chuẩn hóa trước khi so khớp.

\subsubsection{So khớp một hoặc nhiều role}

Rule match khi ít nhất một role của context nằm trong danh sách role được cấu hình.

\subsubsection{Quan hệ giữa role targeting và user targeting}

Allowlist có precedence cao hơn role targeting, vì một người dùng được chỉ định trực tiếp nên được quyết định trước nhóm vai trò rộng hơn.

\subsubsection{Xử lý khi context không chứa role}

Nếu thiếu role, các rule role targeting bị bỏ qua và engine tiếp tục xét lớp rule tiếp theo.

\subsection{Rollout theo phần trăm người dùng}

Percentage rollout dùng stable hashing để phân bổ một targeting key vào bucket cố định từ 0 đến 99.

\subsubsection{Mục tiêu của percentage rollout}

Mục tiêu là phát hành dần tính năng cho một tỷ lệ người dùng ổn định, tránh thay đổi kết quả ngẫu nhiên giữa các request.

\subsubsection{Stable hashing}

Hệ thống băm tổ hợp project/environment/flag/targeting key để cùng một người dùng giữ kết quả ổn định trong phạm vi flag và environment.

\begin{lstlisting}[
    caption={Tạo giá trị băm ổn định bằng SHA-256},
    label={lst:stable-hashing}
]
const hashInput = `${projectKey}:${environmentKey}:${flagKey}:${targetingKey}`;
const digest = sha256(hashInput);
\end{lstlisting}

\subsubsection{Thành phần đầu vào của hàm băm}

Đầu vào băm bao gồm các khóa cấu hình và targeting key. Không sử dụng thời gian hiện tại hoặc số ngẫu nhiên để đảm bảo tính xác định.

\subsubsection{Chuyển đổi kết quả băm thành bucket}

Một phần digest được chuyển thành số nguyên rồi lấy modulo 100 để tạo bucket phần trăm.

\begin{lstlisting}[
    caption={Chuyển đổi SHA-256 digest thành bucket phần trăm},
    label={lst:hash-to-percentage-bucket}
]
const bucket = parseInt(digest.slice(0, 8), 16) % 100;
const enabled = bucket < percentage;
\end{lstlisting}

\subsubsection{Persistent assignment}

Nếu input không đổi, bucket không đổi; vì vậy người dùng không bị nhảy qua lại giữa enabled và disabled.

\subsubsection{Ảnh hưởng khi thay đổi targeting key}

Thay đổi targeting key làm thay đổi bucket, do đó ứng dụng cần chọn khóa ổn định ngay từ đầu.

\subsection{Thực thi rule theo độ ưu tiên}

Rule chỉ được thực thi sau khi các điều kiện hệ thống đã cho phép flag tham gia evaluation.

\subsubsection{Sắp xếp rule}

Trong từng loại rule, các rule enabled được xét theo priority tăng dần.

\subsubsection{Type precedence và first match}

Engine áp dụng type precedence cố định: allowlist, role targeting, percentage rollout. Khi một rule match, engine trả kết quả ngay.

\subsubsection{Rule không hợp lệ hoặc không đủ context}

Rule sai tham số hoặc thiếu dữ liệu context bị bỏ qua thay vì làm bật tính năng ngoài ý muốn.

\subsection{Variant, EvaluationResult và reason code}

Kết quả evaluation phải vừa đủ cho ứng dụng dùng, Dashboard/debug hiểu được, nhưng không phơi bày toàn bộ cấu hình quản trị.

\subsubsection{Variant và response shape}

\begin{lstlisting}[
    caption={Cấu trúc EvaluationResult},
    label={lst:evaluation-result-structure}
]
{
  "enabled": true,
  "variant": "on",
  "reason": "USER_ALLOWLIST",
  "projectKey": "ecommerce-demo",
  "flagKey": "new-checkout"
}
\end{lstlisting}

\subsubsection{Tập reason code}

\begin{lstlisting}[
    caption={Tập reason code của Evaluation API},
    label={lst:evaluation-reason-codes}
]
NOT_FOUND
FLAG_ARCHIVED
FLAG_DISABLED
GROUP_KILL_SWITCH
KILL_SWITCH
GLOBAL_ON
USER_ALLOWLIST
ROLE_MATCH
PERCENTAGE_ROLLOUT
DEFAULT_OFF
INVALID_CONTEXT
ERROR
\end{lstlisting}

\begin{table}[H]
    \centering
    \caption{Reason code cho trạng thái hệ thống}
    \label{tab:system-reason-codes}
    \begin{tabular}{p{0.36\textwidth} p{0.54\textwidth}}
        \toprule
        \textbf{Reason} & \textbf{Ý nghĩa} \\
        \midrule
        \texttt{NOT\_FOUND} & Không tìm thấy project hoặc flag. \\
        \texttt{FLAG\_ARCHIVED} & Flag không còn phục vụ runtime. \\
        \texttt{FLAG\_DISABLED} & Config theo environment đang disabled. \\
        \texttt{GROUP\_KILL\_SWITCH} & Group kill switch đang bật. \\
        \texttt{FLAG\_KILL\_SWITCH} & Kill switch riêng của flag đang bật. \\
        \texttt{ERROR} & Lỗi không mong muốn, trả fail-closed. \\
        \bottomrule
    \end{tabular}
\end{table}

\begin{table}[H]
    \centering
    \caption{Quan hệ giữa mã lý do và matched rule}
    \label{tab:reason-matched-rule}
    \begin{tabular}{p{0.36\textwidth} p{0.54\textwidth}}
        \toprule
        \textbf{Reason} & \textbf{Matched rule} \\
        \midrule
        \texttt{USER\_TARGET\_MATCH} & Rule allowlist match \texttt{userId}. \\
        \texttt{ROLE\_TARGET\_MATCH} & Rule role targeting match role của context. \\
        \texttt{PERCENTAGE\_ROLLOUT\_MATCH} & Bucket stable hashing nằm trong tỷ lệ rollout. \\
        \texttt{DEFAULT\_OFF} & Không rule nào match. \\
        \bottomrule
    \end{tabular}
\end{table}

\subsubsection{Metadata phục vụ truy vết}

Response có thể kèm metadata như matched rule ID hoặc cache status khi cần debug, nhưng metadata này không bắt buộc đối với ứng dụng khách.

\subsection{Default off và error fallback}

Nguyên tắc chung là fail-closed: khi thiếu thông tin hoặc gặp lỗi, hệ thống ưu tiên tắt tính năng mới.

\subsubsection{Nguyên tắc fail-safe và fail-closed}

Fail-closed giúp lỗi hạ tầng không trở thành hành vi mở tính năng ngoài kế hoạch.

\begin{lstlisting}[
    caption={Kết quả fallback mặc định},
    label={lst:default-fallback-result}
]
{
  "enabled": false,
  "variant": "off",
  "reason": "ERROR"
}
\end{lstlisting}

\subsubsection{Flag không tồn tại}

Không tìm thấy project hoặc flag trả \texttt{enabled=false} và \texttt{reason=NOT\_FOUND}.

\subsubsection{Context không đầy đủ}

Context thiếu dữ liệu chỉ làm rule tương ứng không match; nếu không còn rule nào match thì kết quả cuối là \texttt{DEFAULT\_OFF} hoặc \texttt{INVALID\_CONTEXT} tùy loại lỗi.

\subsubsection{Lỗi database hoặc cache}

Khi không tải được snapshot, backend không suy đoán kết quả mà trả fallback tắt.

\subsubsection{Lỗi nội bộ của Evaluation Engine}

\begin{lstlisting}[
    caption={Xử lý lỗi nội bộ của Evaluation Engine},
    label={lst:evaluation-engine-error-handling}
]
try {
  return evaluateWithSnapshot(snapshot, context);
} catch (error) {
  return disabled("ERROR");
}
\end{lstlisting}

\subsection{Khả năng mở rộng Rule Evaluation Engine}

Engine được thiết kế theo hướng thêm loại rule mới mà không thay đổi contract response. Các hướng mở rộng phù hợp gồm targeting theo thuộc tính, organization, thời gian, vị trí và phân phối nhiều variant cho A/B testing.
